\section{The resulting lower bound}

We have obtained the chain estimates needed by the deficit count. We now
apply the earlier word reduction and give the parameters together. This
is the same capacity calculation illustrated in \eqref{eq:counting-example},
with the forest theorem accounting for the components and their endpoints.

For $p\ge0$ and $t>0$, define
\begin{equation}\label{eq:bound}
 \mathcal B_K(p,t)=F_3(K)+K-3+
 \ceil{\frac{2(K-2)!-2(K-2)-4p/t}
 {K^2-4K+1-2(K-1)/t}}.
\end{equation}
In each application below the fraction is positive. The preceding
sections have supplied $p$ and $t$ by bounding the total length of a path
of chains. The following proposition states precisely how that estimate
enters the word bound.

\begin{proposition}[Chain-to-word transfer]\label{prop:transfer}
Let $n=K-1\ge7$. Suppose $p,t$ are real numbers satisfying
\[
 0\le p\le t,\qquad t>0,\qquad t(n-4)\ge4p+6.
\]
If $\CP(p,t)$ holds, then $S(K)\ge\mathcal B_K(p,t)$.
\end{proposition}


The proof combines the scalar deficit estimate with the hereditary
forest theorem. For reference, it is included with the scalar argument
in Appendices~\ref{app:scalar} and~\ref{app:forest}. Its quantifier over all genuine segments
is what permits the estimate to be used on the paths produced by the
reduction from an arbitrary word.

The two return orders in the next theorem are $R_0$ from
\eqref{eq:R0} and $R_1=\ord(M_1)$ from \eqref{eq:efficient-map}.
Their exact values can make the bound stronger at particular sizes.
Replacing them by their common upper bound $n^2/4$ gives a simpler
formula consisting of rational functions on the three residue classes.

\begin{theorem}\label{thm:main}
For every $K\ge8$, we have $S(K)\ge\mathcal B_K(p,t)$ with the following
parameters. If $K\ge13$, write $n=K-1=3m+r$, where $0\le r\le2$, and use
\[
\renewcommand{\arraystretch}{1.5}
\begin{array}{c|c|c}
r&p&t\\\hline
0&2m+\dfrac12+\dfrac1{4R_0}&4m-3-\dfrac1{2R_0}\\
1&m(2m+1)&2m^2+2m-2-\dfrac1{R_1}\\
2&(m+1)(2m+1)&4m^2+5m-\dfrac52
\end{array}
\]
For $8\le K\le12$, use
\[
\renewcommand{\arraystretch}{1.2}
\begin{array}{c|ccccc}
K&8&9&10&11&12\\\hline
p&10&15&521/80&21&28\\
t&46/3&287/12&359/40&263/12&195/4.
\end{array}
\]
In the general table, either return order may instead be replaced by
$n^2/4$.
\end{theorem}


\begin{proof}[Proof of Theorem~\ref{thm:main}]
The chain inequalities are \eqref{eq:CP1}, \eqref{eq:CP2}, and
Proposition~\ref{prop:CP0}. At $K=8$, the earlier theorem gives
$\CP(10,10)$, which implies $\CP(10,46/3)$.
The numerical hypotheses of Proposition~\ref{prop:transfer} are verified
in Appendix~\ref{app:numerical}, including all five initial sizes.
That proposition gives the word bound. The counts remain valid when a
return order is replaced by any larger real number, proving the final
assertion as well.
\end{proof}

Write $B(K)$ for the exact-order formula in Theorem~\ref{thm:main}
and $Q(K)$ for the quadratic-order formula. For $8\le K\le12$ these
use the same initial rows. Every occurrence of a difference between
these formulas below is an ordinary difference of integers.

\subsection{The contribution of the return orders}

To isolate the return-order improvements, define $O(K)$ by the same
initial rows for $K<13$ and, for $K\ge13$, by
$O(K)=\mathcal B_K(p,t)$ with the parameters
\[
\begin{array}{c|c|c}
n=K-1 & p & t\\\hline
3m & 2m & 4m-3\\
3m+1 & m(2m+1) & 2m^2+2m-2\\
3m+2 & (m+1)(2m+1) & 4m^2+5m-5/2.
\end{array}
\]
Thus $O$ already includes the collision between two efficient pairs in
Section~\ref{sec:pairs}. The remaining differences come from the
orders $R_0$ and $R_1$.

\begin{proposition}\label{prop:gains}
For every $K\ge13$ with $K\equiv1,2\pmod3$, we have
$B(K)>O(K)$. Along these two residue
classes, respectively,
\[
 \liminf_{K\to\infty}
 \frac{B(K)-O(K)}{(K-9)!}
 \ge\frac92,
 \qquad
 \liminf_{K\to\infty}
 \frac{B(K)-O(K)}{(K-11)!}
 \ge324.
\]
Replacing $B$ by $Q$, both normalized differences converge to the
displayed constants.
\end{proposition}

\begin{proof}
Before taking ceilings, the exact-order improvements in the two cases are
asymptotic to
\[
 \frac{9n!}{8R_0n^6},\qquad \frac{81n!}{R_1n^8},
\]
respectively. Appendix~\ref{app:numerical} gives the exact subtractions
and verifies that the real improvement exceeds one from $n=12$ and
$n=13$, respectively. The bound $R_i\le n^2/4$ gives the two lower
limits, with equality in the limit for the quadratic choices.
Taking ceilings changes each difference by less than one.
\end{proof}

The exact orders sometimes give larger gains: for instance,
$R_0=2n/7$ when $7\mid n$, and $R_1=n$ in two residue classes of
$m$. This is why Proposition~\ref{prop:gains} states lower limits
for the strongest formula. At $K=13,14$, its gains are $12$ and
$58$; at $K=16,17$, they are $9\,528$ and $16\,920$.

\subsection{Comparison with the all-size reference bound}

Halfar~\cite{halfar} proves the following all-size formula, which he
attributes to Fritsch:
\begin{equation}\label{eq:H}
 H(K)=F_3(K)+K-3+
 \ceil{\frac{2((K-2)!-(K-2))}{K^2-4K+1}},\qquad K\ge7.
\end{equation}
This is Theorem B of the proof package announced on 6 October
2026~\cite{halfar-announcement}. His Theorem C gives stronger bounds
from searched parameters at individual sizes. The comparison below
is with the displayed formula $H$, not those finite bounds.

The earlier chain theorem already improves the all-size formula
\eqref{eq:H} at a larger order than the refinements above. We record
the comparison to distinguish these two contributions.

\begin{corollary}\label{cor:comparison}
For either $A=B$ or $A=Q$, we have
\[
 A(K)=F_3(K)+(2+o(1))(K-4)!,
\]
and, as $K\to\infty$ through the indicated residue class,
\[
 A(K)-H(K)=
 \begin{cases}
 (3+o(1))(K-6)!,&K\equiv1\pmod3,\\
 (18+o(1))(K-7)!,&K\equiv2\pmod3,\\
 (9+o(1))(K-7)!,&K\equiv0\pmod3.
 \end{cases}
\]
\end{corollary}

\begin{proof}
For both parameter choices, $p/t$ is bounded and $t\ge1$ for all
sufficiently large $K$. The numerator inside the ceiling in
\eqref{eq:bound} is therefore $2(K-2)!+O(K)$ and its denominator is
$K^2(1+O(1/K))$. Since $(K-2)!/K^2\sim(K-4)!$, the claimed leading
term follows. The additive term $K-3$ and the rounding error are
$o((K-4)!)$.
With $\kappa=n^2-2n-2$, the difference of the real expressions
inside the ceilings in \eqref{eq:bound} and \eqref{eq:H} is exactly
\[
 \frac{4[n!-n(n-1)-p\kappa]}{t\kappa(\kappa-2n/t)}.
\]
In the three cases $n\equiv0,1,2\pmod3$, the parameter $t$ is
asymptotic to $4n/3$, $2n^2/9$, and $4n^2/9$, respectively.
Substitution gives the three constants. Rounding changes each
difference by less than one.
\end{proof}

\begin{corollary}\label{cor:word-asymptotic}
For every real $\varepsilon>0$, there is $K_0$ such that every
superpermutation on $K\ge K_0$ symbols has length at least
\[
 F_3(K)+(2-\varepsilon)(K-4)!.
\]
\end{corollary}

\begin{proof}
The normalized correction of $B(K)$ tends to $2$, and every covering
word has length at least $B(K)$ by Theorem~\ref{thm:main}.
\end{proof}

The conclusion concerns a lower bound, not an asymptotic formula for
$S(K)$. In particular, the present argument does not produce a positive
coefficient of $(K-3)!$. It strengthens the count of the cheapest local
patterns while retaining the same segment-to-word reduction.
